\documentclass[10pt]{article}
\usepackage[T1]{fontenc}
\usepackage[utf8]{inputenc}
\usepackage[french]{babel}
\usepackage{amsmath,amssymb,mathrsfs,mathpazo}
\usepackage{geometry}
\geometry{a4paper,margin=2cm}
\usepackage{graphicx,array,tabularx,tabularray,colortbl,multirow,multicol,booktabs,enumitem,pifont,ulem,fancybox,diagbox,cancel}
\usepackage{pstricks,pstricks-add,pst-plot,pst-tree,pst-node,pst-text,pst-func,pst-math,pst-bspline,pst-3dplot,pst-eucl,pst-solides3d}
\let\pscalculate\undefined
\usepackage{pst-bezier}
\usepackage{tikz,pgfplots,tkz-tab}
\usetikzlibrary{arrows,arrows.meta,patterns,calc,positioning,trees,decorations.markings,shapes.geometric,shapes.misc,shapes.arrows,intersections,decorations.pathmorphing,decorations.pathreplacing,backgrounds,shadows,babel}
\pgfplotsset{compat=1.18}
\usepgfplotslibrary{fillbetween}
\usepackage[np]{numprint}
\usepackage{listings,makecell,pas-tableur,twoopt}
\let\barre\undefined
\usepackage[pstricks]{tablvar}
\definecolor{darkgreen}{RGB}{0,100,0}
\providecommand{\red}{\color{red}}
\providecommand{\blue}{\color{blue}}
\providecommand{\green}{\color{green}}
\providecommand{\black}{\color{black}}
\providecommand{\white}{\color{white}}
\providecommand{\gray}{\color{gray}}
\providecommand{\cyan}{\color{cyan}}
\providecommand{\magenta}{\color{magenta}}
\providecommand{\yellow}{\color{yellow}}
\providecommand{\violet}{\color{violet}}
\providecommand{\orange}{\color{orange}}
\providecommand{\euro}{\texteuro}
\providecommand{\EUR}{\texteuro}
\providecommand{\ieme}{\textsuperscript{e}}
\providecommand{\degres}{\ensuremath{^\circ}}
\providecommand{\degree}{\ensuremath{^\circ}}
\providecommand{\vv}[1]{\overrightarrow{#1}}
\providecommand{\Coord}[1]{#1}
\providecommand{\decosix}{}
\providecommand{\starredbullet}{\ensuremath{\bullet}}
\pagestyle{empty}
\setlength{\parindent}{0pt}
\usepackage{adjustbox}\DeclareUnicodeCharacter{2236}{:}\DeclareUnicodeCharacter{2190}{\ensuremath{\leftarrow}}
\providecommand{\R}{\mathbb{R}}
\providecommand{\N}{\mathbb{N}}
\providecommand{\Z}{\mathbb{Z}}
\providecommand{\Q}{\mathbb{Q}}
\providecommand{\D}{\mathbb{D}}
\providecommand{\C}{\mathbb{C}}
\providecommand{\e}{{\mathrm{e}}}
\providecommand{\vect}[1]{\overrightarrow{#1}}
\providecommand{\vectt}[1]{\overrightarrow{\mathrm{#1}}}
\providecommand{\vvt}[1]{\overrightarrow{\mathrm{#1}}}
\providecommand{\np}[2][]{{#2~#1}}
\providecommand{\vv}[1]{{\overrightarrow{#1}}}
\renewcommand{\d}{{\mathrm{d}}}
\providecommand{\degres}{{^\circ}}
\providecommand{\degre}{{^\circ}}
\providecommand{\blue}{\color{blue}}
\providecommand{\red}{\color{red}}
\providecommand{\enskip}{\quad}
\providecommand{\llbracket}{\lbrack\!\lbrack}
\providecommand{\rrbracket}{\rbrack\!\rbrack}
\providecommand{\euro}{{\text{€}}}
\providecommand{\no}{{\text{n°}}}
\providecommand{\interval}[2]{{[#1;#2]}}
\providecommand{\intervFO}[2]{{[#1;#2[}}
\providecommand{\dotfill}{\ldots}

\providecommand{\ds}{\displaystyle}
\providecommand{\textding}[1]{#1}
\providecommand{\Oij}{\(\mathrm{O};\vec{\imath},\vec{\jmath}\)}
\providecommand{\Oijk}{\(\mathrm{O};\vec{\imath},\vec{\jmath},\vec{k}\)}
\providecommand{\Ouv}{\(\mathrm{O};\vec{u},\vec{v}\)}
\providecommand{\intervFF}[2]{\ensuremath{[#1;#2]}}
\renewcommand{\euro}{\ifmmode\text{\texteuro}\else\texteuro\fi}
\DeclareUnicodeCharacter{00B0}{\ifmmode^{\circ}\else\textdegree\fi}
\DeclareUnicodeCharacter{20AC}{\ifmmode\text{\texteuro}\else\texteuro\fi}
% Énoncés : Éducation nationale. Transcription et corrigés : APMEP, Denis Vergès et contributeurs. Découpage et classement : CoopMaths, Sébastien Lozano, Éric Elter et contributeurs. Adaptation de la présentation : Maths Explorer.
% CC BY-SA : https://coopmaths.fr/www/about/
% https://coopmaths.fr/alea/static/bac/2021/tex/bac_2021_06_etrangers_9.tex
\begin{document}
\noindent\textit{APMEP / CoopMaths. CC BY-SA. Présentation adaptée par Maths Explorer.}\par\medskip
\bigskip

\textbf{Partie A :}

\medskip

Dans un pays, une maladie touche la population avec une probabilité de $0,05$. 

On possède un test de dépistage de cette maladie.

On considère un échantillon de $n$ personnes ($n \geqslant 20$) prises au hasard dans la population assimilé à un tirage avec remise.

On teste l'échantillon suivant cette méthode : on mélange le sang de ces $n$ individus, on teste le mélange. 

Si le test est positif, on effectue une analyse individuelle de chaque personne.

Soit $X_n$ la variable aléatoire qui donne le nombre d'analyses effectuées.

\medskip

\begin{enumerate}
\item Montrer $X_n$ prend les valeurs 1 et $(n + 1)$. 
\item Prouver que $P\left(X_n = 1\right) = 0,95^n$.

Établir la loi de $X_n$ en recopiant sur la copie et en complétant le tableau suivant: 

\begin{center}
\begin{tabularx}{0.6\linewidth}{|*{3}{>{\centering \arraybackslash}X|}}\hline
$x_i$& 1 &$n + 1$\\ \hline
$P\left(X_n = x_i\right)$&&\\ \hline
\end{tabularx}
\end{center}

\item Que représente l'espérance de $X_n$ dans le cadre de l'expérience ?

Montrer que $E\left(X_n\right) =n + 1 - n \times  0,95^n$. 
\end{enumerate}

\bigskip

\textbf{Partie B :}

\medskip

\begin{enumerate}
\item On considère la fonction $f$ définie sur $[20~;~ +\infty[$ par $f(x) = \ln (x) + x \ln (0,95)$.

Montrer que $f$ est décroissante sur $[20~;~ +\infty[$.
\item On rappelle que $\displaystyle\lim_{x \to + \infty} \dfrac{\ln x}{x} = 0$. Montrer que $\displaystyle\lim_{x \to + \infty} f(x) = - \infty$.
\item Montrer que $f(x) = 0$ admet une unique solution $a$ sur $[20~;~ +\infty[$. 

Donner un encadrement à 0,1 près de cette solution.
\item En déduire le signe de $f$ sur $[20~;~ +\infty[$. 
\end{enumerate}

\bigskip

\textbf{Partie C :}

\medskip

On cherche à comparer deux types de dépistages. 

La première méthode est décrite dans la partie A, la seconde, plus classique, consiste à tester tous les individus. 

La première méthode permet de diminuer le nombre d'analyses dès que $E\left(X_n\right)  < n$.

En utilisant la partie B, montrer que la première méthode diminue le nombre d'analyses pour des échantillons comportant $87$ personnes maximum.

\bigskip

\end{document}
